\chapter{Organização Documental da Alern}
\label{cap:DocAlern}

Conforme apontado no Capítulo \ref{cap:Introducao}, a produção de dados em formato de texto (quer estruturado ou não) é uma característica intrínseca da dos processos integrantes da atividade legislativa da \ac{Alern}. Tal produção é forçosamente acompanhada de uma atividade subjacente de tratamento, armazenamento e manutenção desses dados, tanto para fins de registros de uso interno, como para fins de publicização das informações ao escrutínio público. Neste capítulo apontamos precisamente as duas frentes mencionadas, a saber: as ações tomadas no intuito de disponibilização dos dados ao cidadão, assim como a manutenção das informações internas da própria Casa Legislativa.

\section{Iniciativas de Promoção de Transparência}
\label{sec:DocAlern.1}

A Assembleia Legislativa do Rio Grande do Norte (\ac{Alern}) frequentemente se  apresenta como uma instituição que mantém um compromisso crescente com a transparência e o acesso à informação, apontando diversas iniciativas e projetos voltados a ampliar a interação e o controle social sobre suas atividades legislativas e administrativas. Nesse sentido, a Casa Legislativa tem buscado desenvolver cursos de ação e medidas com foco na busca por ``ampliar os índices de acesso à informação e transparência'' e na garantia do respeito aos ``critérios de Transparência e Fiscalização nacionalmente consolidados por órgãos de Controle e Gestão como Tribunais de Contas e Controladoria-Geral da União''. \cite{ALERN:1}

Tal qual outros órgãos e instituições públicas dos três poderes, nas esferas federal, estadual e municipal, a forma oficial de a \ac{Alern} levar ao conhecimento do cidadão as matérias referentes a sua atividade legislativa se dá por meio de edições regulares do \textbf{Diário Oficial Eletrônico} \cite{ALERN:4}. O Diário Oficial inclui todos os atos administrativos da Casa, bem como toda a publicação de informações de interesse público referentes à sua atuação. Assuntos adicionais, coberturas de eventos, sumarização de discussões e atuações da atividade parlamentar também são publicadas em forma de notícias na página oficial do órgão.

Em um tom mais voltado à prestação de contas, são disponibilizadas, desde 2017, no \textbf{Portal da Transparência} \cite{ALERN:3}, informações de gestão administrativa e financeira, tratando de questões orçamentárias, receitas e despesas, folha de pagamento e diárias dos servidores e contratos; além de assuntos abrangidos tanto pela Lei de Acesso à Informação quanto pela Lei de Responsabilidade Fiscal, servindo como instrumento de participação e fiscalização popular. Porém, particularmente a partir do início de 2024, ocorre uma maior publicidade e transparência no processo legislativo da \ac{Alern} em diversas áreas, quando, além da página de transparência já antes oferecida ao escrutínio público, foi acrescentada também a página de \textbf{Transparência Legislativa} \cite{ALERN:5}, visando garantir maior acesso da população ``à rotina do trabalho parlamentar em tempo real'', permitindo-lhes ``acompanhar a agenda legislativa, incluindo sessões plenárias, reunião de comissões, audiências públicas e sessões solenes, além da tramitação de leis e projetos do legislativo estadual''. \cite{ALERN:2}

Além disso, também se tornou

 \begin{citacao}
 possível acompanhar o trabalho dos 24 deputados estaduais de forma individualizada. Cada parlamentar possui uma página institucional com informações da respectiva atuação legislativa, com projetos apresentados, aprovados, participação em comissões temáticas e notícias mais recentes. \cite{ALERN:2}
 \end{citacao}

Juntamente com a divulgação ativa para consumo passivo da população, foi feita a implementação de ferramentas que permitem a inversão desses papéis, garantindo a possibilidade de uma busca ativa do cidadão sobre um conteúdo estático, através de busca textual avançada sobre todo o conteúdo disponível, ``com adição de filtros, como período, região, parlamentar propositor, assunto, número do projeto, abrangência, e outros itens''. \cite{ALERN:2} Embora seja um aprimoramento relevante para as opções que a plataforma oferece ao cidadão, o conjunto de critérios possíveis de se aplicar na busca são limitantes, visto que só é possível buscar por termos exatos, não dando conta, por exemplo, de sinônimos ou variações na grafia. Como consequência, ocorre uma subidentificação de conteúdo relevante, incorrendo em omissão de informações com potencial de interesse, diminuindo a abrangência da ferramenta de busca. Diante das limitações existentes, cabe ao usuário saber de antemão quais os termos exatos um parlamentar usou para expressar uma ideia -- situação em que uma busca por similaridade semântica dos termos (ver Seção \ref{sec:Teorico.3}) poderia aprimorar os resultados, não se limitando a palavras específicas.

Através da plataforma \textbf{Legis Vídeos} \cite{ALERN:6}, a Assembleia Legislativa potiguar garante à população o acesso às gravações dos eventos legislativos por si realizados, remontando a 2015, incluindo sobretudo registros das sessões e reuniões dos parlamentares e das comissões. Em termos de informações em formato de texto, de interesse maior deste estudo, cabe apontar que na mesma plataforma de acesso aos registros videográficos dos eventos, quando se trata de sessões parlamentares, são disponibilizadas as atas em formato textual com o resumo das suas respectivas reuniões. Adicionalmente a isso, são disponibilizadas as transcrições do áudio capturado na gravação das sessões, as quais são produzidas de forma automatizada, através de ferramentas de conversão de áudio para texto que se utilizam de modelos de inteligência artificial. \cite{ALERN:3}

Essas últimas iniciativas de publicidade, em particular, possuem um potencial informativo extremamente expressivo, em termos de fazer saber sobre a atividade parlamentar de forma mais abrangente. Disponibilizar o produto final da atuação dos Deputados -- na forma do texto das leis aprovadas, sancionadas e promulgadas -- é parte da atividade obrigatória da Casa Legislativa, sem a qual o exercício da cidadania através do conhecimento da regulação social é impraticável; pelo que mais que se justifica a publicidade do processo legislativo por meio de ferramentas como o Diário Oficial. Não obstante, para o esclarecimento do cidadão, tão importante quanto o conhecer das leis em si é acompanhar os processos de discussão que culminaram no texto final de uma peça normativa em si. Através da análise do discurso, do entendimento das interações entre os parlamentares, da compreensão das ideias propostas durante as sessões plenárias, torna-se mais sólido o processo de acompanhamento do mandato do parlamentar.

\section{Gerenciamento de Informação Interna}
\label{sec:DocAlern.2}

Quando da escrita deste estudo, todo o registro do processo de redação, proposição, discussão, edição, votação e aprovação das leis (chamado processo legislativo) é realizado através do \textbf{e-Legis}, tanto por parte dos parlamentares da Assembleia e quanto pela estrutura de servidores que dão suporte à atividade parlamentar. Trata-se de uma premiada ferramenta desenvolvida pelas equipes da própria Casa Legislativa com o intuito de dinamizar a rotina do órgão e manter um registro acurado dos processos como um todo. Correntemente, as informações disponibilizadas pela mencionada página de Transparência Legislativa são um recorte dos dados gerados durante os trâmites do processo legislativo, conforme conduzido pela \ac{Alern}.

Como resultado da utilização da plataforma, são gerados registros de cada uma das versões dos documentos que, ao final da sequência de ações tomadas pelos Deputados, tornar-se-á o texto final de uma lei ou peça normativa semelhante. Juntamente com os dados textuais, um conjunto de metadados -- gerados e coletados no transcorrer de cada uma das partes da tramitação -- é armazenado de forma persistente, o que acaba por criar um corpus informativo capaz de descrever todo o fluxo do processo legislativo no que toca a um projeto/lei. Atualmente, há estudos sendo feitos na análise dos dados e metadados gerados no processo legislativo, com o fim de identificar padrões e semelhanças entre novas possíveis proposições e matérias análogas já tramitadas na Casa.

Considerando todo o cabedal de informações geradas em todo o âmbito da \ac{Alern}, tal qual nas organizações de cunho público em geral, a sua utilização, processamento e submissão a ferramentas de pesquisa e busca disponíveis ao indivíduo comum se torna não somente uma possibilidade, mas igualmente uma responsabilidade a ser cumprida, a fim de se alcançar o objetivo da transparência na coisa pública. O desenvolvimento, portanto, de ferramentas que fruam a abundância de dados produzidos e mantidos por instituições como a \ac{Alern} para fomentar o processo de prestação de contas junto à sociedade é uma atividade extremamente relevante e desejável que seja realizada, tanto pelas próprias entidades quanto pela sociedade civil.

Nesse contexto, considerando o desenvolvimento e a disponibilidade de \textit{software} de processamento de dados -- especialmente dados textuais, conforme será abordado no Capítulo \ref{cap:Teorico} -- o desenvolvimento de soluções que permitam a extração, processamento e consulta eficiente das informações se demonstra viável e factível.